% Figure from MATH 451 notes, section 2. Compile with the notes preamble.
\begin{tikzpicture}[scale=1.0]
  \draw[-stealth, thin] (0.4,0.4) -- (6.0,0.4) node[right, font=\scriptsize] {$p$};
  \draw[-stealth, thin] (0.4,0.4) -- (0.4,6.0) node[above, font=\scriptsize] {$q$};
  \foreach \i in {1,...,5} {
    \draw[thin] (\i,0.47) -- (\i,0.33) node[below, font=\scriptsize] {$\i$};
    \draw[thin] (0.47,\i) -- (0.33,\i) node[left, font=\scriptsize] {$\i$};
  }
  % diagonals p+q = n
  \foreach \n in {2,...,6} {
    \draw[black!30, dashed] ({\n-0.75},0.75) -- (0.75,{\n-0.75});
    \node[font=\scriptsize, black!55] at ({\n-0.88},0.62) {$\n$};
  }
  % order of listing (within each diagonal: increasing p)
  \draw[figR, thin, -stealth] (1,2) -- (2,1);
  \draw[figR, thin, -stealth] (1,3) -- (2.9,1.1);
  \draw[figR, thin, -stealth] (1,4) -- (4,1);
  \draw[figR, thin, -stealth] (1,5) -- (4.93,1.07);
  % non-coprime pairs: skipped
  \foreach \p/\q in {2/2,2/4,3/3,4/2} {
    \draw[black!55, fill=white, thin] (\p,\q) circle (2.4pt);
  }
  % coprime pairs: listed
  \foreach \p/\q in {1/1,1/2,2/1,1/3,3/1,1/4,2/3,3/2,4/1,1/5,5/1} {
    \fill[figB] (\p,\q) circle (2.4pt);
    \node[font=\scriptsize, figB, anchor=south west, inner sep=1.5pt] at (\p,\q) {$\tfrac{\p}{\q}$};
  }
  \node[font=\scriptsize, black!70, anchor=west, align=left] at (3.3,4.7) {\textcolor{figB}{$\bullet$} listed ($p,q$ coprime)\\ $\circ$ skipped (already seen)\\ \textcolor{black!55}{dashed: $p+q = 2,3,\ldots,6$}};
\end{tikzpicture}
